\chapter*{Preface}
\addcontentsline{toc}{chapter}{Preface}

This project grew out of a simple interest in parameter-efficient continual
learning: given that low-rank adaptation already lets a large pretrained
model be fine-tuned cheaply for a single task, what happens when that
fine-tuning has to happen repeatedly, one incremental step after another,
without revisiting the data that came before? The central empirical
question that emerged from this starting point was not whether low-rank
adaptation works for continual learning in isolation, but how it should be
\emph{organised} across a sequence of incremental steps.

Two candidate organisations run through this thesis. The first,
\emph{SimpleAvg}, is deliberately simple: each step gets its own
independently trained, fixed-rank adapter, and their reconstructed dense
updates are combined once, at the end of the sequence, by arithmetic
averaging. It is treated throughout as a
reference point rather than as a proposed contribution in itself. The
second, \emph{RankExt}, is a persistent alternative in which a single
adapter is grown incrementally, one rank increment per step, with earlier
blocks left frozen but still active. One question that motivated much of
the experimental work reported here is how much of SimpleAvg's performance
a persistent, growing-rank strategy can recover, and what
stability--plasticity trade-off accompanies that recovery as auxiliary retention
mechanisms are added.

The resulting work is mainly empirical, comparative, diagnostic, and
ablative in character. It compares the two strategies under a common
protocol, examines where and why one struggles relative to the other, and
tests a number of auxiliary mechanisms and refinements against that
diagnosis. It does not claim to have invented low-rank adaptation, model
averaging, growing-rank continual adaptation, knowledge distillation, or
orthogonality-based regularisation; each of these building blocks is drawn
from existing practice and combined here for the specific purpose of this
comparison.

I hope the chapters that follow are useful not only for the specific
numbers they report, but for the diagnostic approach they take to reading
those numbers.
